\documentclass{article}
\usepackage[utf8]{inputenc}
\usepackage[T1]{fontenc}
\usepackage{geometry}
\geometry{
    a4paper,
    left=2.5cm,
    right=2.5cm,
    top=2.5cm,
    bottom=2.5cm
}
\usepackage{graphicx}
\usepackage{float}
\usepackage{hyperref}
\usepackage[english]{babel}
\usepackage{amsmath}
\usepackage{amssymb}
\usepackage{array}
\usepackage{booktabs}
\usepackage{caption}
\usepackage{subcaption}
\usepackage{enumitem}
\usepackage{fancyhdr}
\usepackage{titlesec}
\usepackage{url}
\usepackage{textcomp}
\usepackage{xcolor}
\usepackage{listings}
\usepackage{multirow}
\usepackage{longtable}
\usepackage{tabularx}
\usepackage{times}

% Set up page style
\pagestyle{fancy}
\fancyhf{}
\rfoot{\thepage}
\renewcommand{\headrulewidth}{0pt}

% Title formatting
\titleformat{\chapter}[display]{\normalfont\bfseries}{\chaptertitle\thechapter}{0pt}{\Huge}
\titlespacing*{\chapter}{0pt}{50pt}{40pt}

% Document info
\title{Template Document}
\author{Academic Support Agent}
\date{\today}

\begin{document}
\maketitle

\maketitle

\newpage

\section*{Abstract}
This section provides a concise summary of the proposed project. It should include the project's objectives, methods, expected outcomes, and potential impact. The abstract is often limited to one page.

\newpage

\section{Specific Aims}
State the specific goals of the project, typically in bullet points. Be clear and concise, focusing on what you intend to accomplish.

\newpage

\section{Research Strategy}
\subsection{Significance}
Explain the importance of the problem or critical barrier to progress in the field that the proposed project addresses.

\subsection{Innovation}
Describe any novel theoretical concepts, approaches, methodologies, instrumentation, or interventions to be developed or used.

\subsection{Approach}
Discuss the overall strategy, methodology, and analyses to be used to accomplish the specific aims. Include how the data will be collected, analyzed, and interpreted.

\newpage

\section{Budget Justification}
Provide a detailed justification for each budget category, explaining the necessity and reasonableness of the costs.

\newpage

\section{Biographical Sketches}
Include biographical sketches of key personnel, following the NIH format. Each biosketch is typically limited to five pages.

\newpage

\section{Facilities and Other Resources}
Describe the facilities and resources available for conducting the proposed research. This section should demonstrate that you have access to the necessary infrastructure.

\newpage

\section{References}
Enter your references manually or use refer to a .bib file in your project.

% Manually input references here. For example:
\begin{itemize}
    \item Author1, A. A., Author2, B. B. (Year). Title of the article. \textit{Journal Name}, Volume(Issue), pages.
    \item Author3, C. C., Author4, D. D. (Year). Title of the book. \textit{Publisher}.
    % Add more references as needed.
\end{itemize}

% Use a bib for references here:
%
%
\end{document}